\section{Object-Oriented Programming (OOP)}
\label{sec:oop}

The fundamental principle of OOP is the encapsulation of data (attributes) and logic (methods) into objects. This structure ensures data integrity by restricting direct external manipulation.

\subsection{Class Definition and Instantiation}
A class serves as a formal blueprint. Per PEP 8, class names follow the \texttt{CapWords} convention. Objects are specific instances created from these blueprints.

\begin{lstlisting}[language=Python]
class Cat:
    """Formal structure of a cat."""
    pass

# Instantiation
garfield = Cat()
\end{lstlisting}

\subsection{Attributes and State Integrity}
Attributes store data within an object. Python distinguishes between Class and Instance variables based on their scope and definition.

\subsubsection{Class vs. Instance Variables}
\begin{itemize}
    \item \textbf{Class Variables:} Defined directly within the class body. They are shared across all instances.
    \item \textbf{Instance Variables:} Defined using \texttt{self}. They are unique to each instance.
\end{itemize}\vspace{1em}

\textbf{\color{red}Critical Warning:} Never modify a class variable through an instance 
(e.g., \texttt{instance.class\_var = value}). This creates a new instance attribute 
that shadows the class variable for that specific object.

Always update class variables using the class name: \texttt{ClassName.class\_var = value}.

\subsubsection{Initialization and Dynamic Attributes}
It is highly recommended to define \textbf{all} instance variables within the \texttt{\_\_init\_\_} method.
\begin{itemize}
    \item \textbf{Avoid 'On-the-fly' attributes:} It makes the object's structure unpredictable and harder to debug.
    \item \textbf{Defined State:} Ensuring all attributes are initialized in the constructor
\end{itemize}

\subsection{Method Attributes and self}
Methods are functions defined within a class. The first parameter is always \texttt{self}, which represents the instance calling the method. Python binds the instance to this parameter automatically during the call.

\begin{lstlisting}[language=Python]
    account1.deposit(100)
    # In background will this executed
    Class_name.deposit(self=account1, amount=100) 
\end{lstlisting}

\subsection{Data Abstraction and Encapsulation}
Encapsulation hides implementation details, exposing only necessary information through defined interfaces. Python uses naming conventions to signal the intended visibility of attributes:

\resizebox{\columnwidth}{!}{
\begin{tabular}{lll}
\hline
\textbf{Category} & \textbf{Convention} & \textbf{Description} \\ \hline
Public & \texttt{name} & Accessible from anywhere. \\
Protected & \texttt{\_name} & Undesirable to access externally; used for inheritance. \\
Private & \texttt{\_\_name} & Hidden via name mangling; not directly accessible. \\ \hline
\end{tabular}}\vspace{1em}

A private variable is \texttt{my\_object.\_MyClass\_\_private\_var}

\subsection{Properties}
The property object is created using the signature: \\
\texttt{var = property(fget=None, fset=None, fdel=None, doc=None)}

\lstinputlisting[language=Python]{code/setter_and_getters.py}

\subsubsection{Decorators}

The same could be achieved using decorators:

\lstinputlisting[language=Python]{code/setter_and_getters_decorators.py}

\section{Magic Methods and Custom Behavior}

\subsection{Object Representation}
Magic methods (or dunder methods) allow custom classes to interact with built-in Python functions and operators.
\begin{itemize}
    \item \texttt{\_\_str\_\_(self)}: Returns a user-friendly string representation (used by \texttt{print()}).
    \item \texttt{\_\_repr\_\_(self)}: Returns an unambiguous string representation, ideally used for debugging or recreating the object.
\end{itemize}

\begin{lstlisting}[language=Python]
class Book:
    def __init__(self, title, author):
        self.title = title
        self.author = author

    def __str__(self):
        return f"'{self.title}' by {self.author}"

    def __repr__(self):
        return f"Book(title='{self.title}', author='{self.author}')"
\end{lstlisting}

\subsection{Emulating Container Types}
Classes can behave like sequences (lists) or mappings (dictionaries) by implementing specific methods:
\begin{itemize}
    \item \texttt{\_\_len\_\_(self)}: Returns the number of items.
    \item \texttt{\_\_getitem\_\_(self, key)}: Defines behavior for indexing (\texttt{obj[key]}).
    \item \texttt{\_\_setitem\_\_(self, key, value)}: Defines behavior for item assignment (\texttt{obj[key] = value}).
\end{itemize}

\begin{lstlisting}[language=Python]
class Inventory:
    def __init__(self):
        self._items = {}

    def __setitem__(self, key, value):
        self._items[key] = value

    def __getitem__(self, key):
        return self._items.get(key, "Item not found")

    def __len__(self):
        return len(self._items)
\end{lstlisting}

\subsection{Context Managers}
Context managers ensure resources (like files or network connections) are properly managed using the \texttt{with} statement. This requires the implementation of two methods:
\begin{itemize}
    \item \texttt{\_\_enter\_\_(self)}: Executed at the start of the \texttt{with} block. It returns the resource.
    \item \texttt{\_\_exit\_\_(self, exc\_type, exc\_val, tb)}: Executed at the end, even if an exception occurs.
\end{itemize}

\begin{lstlisting}[language=Python]
class FileManager:
    def __init__(self, filename, mode):
        self.filename = filename
        self.mode = mode

    def __enter__(self):
        self.file = open(self.filename, self.mode)
        return self.file

    def __exit__(self, exc_type, exc_val, traceback):
        self.file.close()

with FileManager("data.txt", "w") as f:
    f.write("Structured data")
\end{lstlisting}

\subsection{Reflected Arithmetic Methods (The 'r' Prefix)}
Python distinguishes between an operation where an object is on the left-hand side (LHS) and one where it is on the right-hand side (RHS). This is crucial for interoperability with built-in types (like \texttt{int} or \texttt{float}).

\begin{itemize}
    \item \textbf{\texttt{\_\_add\_\_(self, other)}}: Called for \texttt{self + other}.
    \item \textbf{\texttt{\_\_radd\_\_(self, other)}}: Called for \texttt{other + self} only if the LHS object does not define \texttt{\_\_add\_\_} for the specific type of \texttt{self}.
\end{itemize}

\begin{lstlisting}[language=Python]
class Price:
    def __init__(self, amount):
        self.amount = amount

    def __add__(self, other):
        # Handles: Price(10) + 5
        val = other.amount if isinstance(other, Price) else other
        return Price(self.amount + val)

    def __radd__(self, other):
        # Handles: 5 + Price(10)
        # Re-uses __add__ since addition is commutative
        return self.__add__(other)
\end{lstlisting}

\subsection{The \texttt{NotImplemented} Singleton}
In Python, \texttt{NotImplemented} è un oggetto speciale (un singleton) restituito dai metodi binari (come \texttt{\_\_add\_\_}, \texttt{\_\_sub\_\_}, \texttt{\_\_eq\_\_}, ecc.) per indicare che l'operazione non è implementata per il tipo di dato dell'altro operando.

\textbf{Importante:} Non deve essere confuso con \texttt{NotImplementedError}, che è un'eccezione. \texttt{NotImplemented} è un valore di ritorno che segnala a Python di tentare una strategia alternativa.
